\section{Quotient decompositions, linear duality and sharp transitions}
\label{sec:r74-quotient-memory}

For groups whose subgroups are normal, the probability optimization in the
normal form has an exact linear reduction. This includes every finite abelian
group and allows arbitrary probability tables and terminal outcome sets.
Neither a special phase choice nor a positive-amplitude trigonometric ansatz
is imposed. The reduction uses normal stabilizers; it is not an unrestricted
claim that equivariant nonnegative factorization is linear programming.

\begin{lemma}[Absorption of within-orbit preparation]
\label{lem:r74-normal-orbits}
Suppose the stabilizers in a finite $G$-set are normal. Every probability
factorization on it can be replaced, without increasing its size, by one
having at most one orbit $H\backslash G$ for each normal subgroup $H$ and
whose initial row is concentrated at the identity coset within each orbit.
Only the mixture between different orbits is required.
\end{lemma}
\begin{proof}
On an orbit of positive initial mass, normalize its local preparation and
write its response as
\[
 q(g,c)=\sum_{Hx\in H\backslash G}\lambda(Hx)D(Hxg,c).
\]
For $h\in H$, normality gives $Hxhg=Hxg$, so $q(hg,c)=q(g,c)$.
Thus $q$ itself is a probability decoder on $H\backslash G$. Replacing
the preparation by the identity coset and the decoder by $q$ leaves every
command-word response unchanged. Orbits with the same $H$ can be merged
by averaging their decoders with their initial masses as weights.
Delete zero-mass orbits. These operations weakly decrease the size.
\end{proof}

Assume henceforth, until the Fourier specialization, that every subgroup of
$G$ is normal. Such groups are called Dedekind groups. For a nonempty set
$\mathcal H$ of distinct subgroups, put
$c(\mathcal H)=\sum_{H\in\mathcal H}[G:H]$.

\begin{theorem}[Exact quotient linear programs]
\label{thm:r74-quotient-lp}
The invariant $\kappa_G(b)$ is the minimum $c(\mathcal H)$ for which there
are $p_H\ge0$ and functions
$f_H:H\backslash G\times C\to\R_{\ge0}$ satisfying
\begin{equation}\label{eq:r74-quotient-lp}
 \sum_cf_H(q,c)=p_H\quad(q\in H\backslash G),\qquad
 \sum_Hp_H=1,\qquad
 \sum_H f_H(Hg,c)=b(g,c)\quad(g\in G,c\in C).
\end{equation}
The residual cardinality is $N=|G|$, and
$m(b)=\min\{|G|,1+\kappa_G(b)\}$.
For a rational response table a minimum controller can be chosen rational.
More generally, if the table lies in a specified real algebraic number field,
the controller can be chosen in that same field.
\end{theorem}
\begin{proof}
Apply Lemma~\ref{lem:r74-normal-orbits} to a minimum probability
factorization. Its orbit contributions give \eqref{eq:r74-quotient-lp}.
Conversely, if $p_H>0$, put $D_H=f_H/p_H$, place mass $p_H$ at its identity
coset, and use the right quotient action. If $p_H=0$, all its entries vanish
and the orbit is discarded. This constructs the whole response with at
most $c(\mathcal H)$ points. The regular orbit $H=\{1\}$ is always
available, and there are finitely many subgroup families.

The stabilizer $H_0$ of the initial residual is normal. If $h\in H_0$
and $g\in G$, moving $h$ past $g$ using normality shows that $h$ fixes
the residual after $g$. It therefore fixes their span and is the identity
in the minimal observable group. Thus $H_0=\{1\}$ and orbit-stabilizer
gives $N=|G|$. The pure-start assertion follows from
Theorem~\ref{thm:r73-memory-formula}.

For each family the displayed feasible set is a bounded closed polyhedron:
$0\le f_H(q,c)\le p_H\le1$. A feasible vertex is obtained by solving a
linear system over the coefficient field. It is rational for rational data
and remains in the given real algebraic field otherwise. The construction
uses only these entries and division by positive $p_H$. If the residual
machine is smaller than the initializer construction, its decoder is the
given table and has the same arithmetic property.
\end{proof}

This theorem optimizes the common permutation action supplied by the inherited
normal form. The relation between symmetry and positive factorizations is
classical \cite{GPT13,KPT10}, and positivity can cost more than a factorization
inside the target linear space \cite{GillisGlineur12}. The additional step
here is the absorption of preparation on normal quotient orbits, which
removes the bilinear unknowns without changing the counted resource.

\subsection{An explicit payoff dual}
Let $\mathcal P_H$ be the polytope of all probability tables on
$H\backslash G$, pulled back to $G$. For a real table $y$ on $G\times C$,
use the unnormalized pairing $\langle y,b\rangle=\sum_{g,c}y(g,c)b(g,c)$
and set
\begin{equation}\label{eq:r74-dual-support}
 h_H(y)=\sum_{q\in H\backslash G}\max_{c\in C}\sum_{g\in q}y(g,c),
 \qquad h_{\mathcal H}(y)=\max_{H\in\mathcal H}h_H(y).
\end{equation}

\begin{theorem}[Payoff certificates and exact contrast endpoints]
\label{thm:r74-quotient-dual}
The family $\mathcal H$ is feasible in \eqref{eq:r74-quotient-lp} if and
only if
\begin{equation}\label{eq:r74-payoff-test}
 \langle y,b\rangle\le h_{\mathcal H}(y)
 \quad\hbox{for every real table }y.
\end{equation}
If infeasible, it has a strict rational witness $y$. A witness of margin
$\Delta=\langle y,b\rangle-h_{\mathcal H}(y)>0$ still excludes that family
for any table $\widetilde b$ with
$\|\widetilde b-b\|_\infty<\Delta/\|y\|_1$.

For the contrast line \eqref{eq:r74-contrast}, put $v=b-\pi$. Its endpoint
for this family is
\begin{equation}\label{eq:r74-threshold-dual}
 \tau_{\mathcal H}
 =\min\left\{1,
  \min_{\langle y,v\rangle=1}
       \bigl(h_{\mathcal H}(y)-\langle y,\pi\rangle\bigr)\right\}.
\end{equation}
Both minima are attained. Equivalently, maximize $t$ in $[0,1]$ subject to
\eqref{eq:r74-quotient-lp} with right side $\pi+tv$.
The endpoint for budget $k$ is the maximum of these linear-program values
over $c(\mathcal H)\le k$. Rational tables give rational endpoints and
optimal controllers; no polynomial bound on the number of families is claimed.
\end{theorem}
\begin{proof}
The feasible response set for a fixed family is exactly
$\mathcal P_{\mathcal H}=\operatorname{conv}\bigcup_H\mathcal P_H$.
Each $\mathcal P_H$ is a product of probability simplexes. Maximizing a
payoff independently on each coset gives its support function $h_H$;
taking a convex hull of the union gives $h_{\mathcal H}$.
Membership proves \eqref{eq:r74-payoff-test}. If $b$ is outside this compact
convex set, let $q$ be a closest point. The first variation along the
segment from $q$ to any point of the set gives
$\langle b-q,x-q\rangle\le0$. With $y=b-q$ this yields a strictly
positive gap of at least $\|b-q\|_2^2$. Approximate $y$ by a rational
table sufficiently closely to preserve that gap; the support function is
a finite maximum of linear functions and is continuous. The perturbation
assertion is the inequality
$|\langle y,\widetilde b-b\rangle|\le\|y\|_1
\|\widetilde b-b\|_\infty$.

The constant table $\pi$ belongs to every $\mathcal P_H$. Along its ray
$\pi+tv$, a support inequality with $\langle y,v\rangle\le0$ imposes
no restriction for $t\ge0$; every other inequality can be normalized by
$\langle y,v\rangle=1$. This gives \eqref{eq:r74-threshold-dual}.
Attainment requires no strict-feasibility assumption. If $v$ lies outside
the linear space parallel to the affine hull of $\mathcal P_{\mathcal H}$,
an affine-hull equation gives an attained inner minimum zero. Otherwise
the bounded polytope has a finite facet description in that affine hull;
the first boundary reached along the ray activates a facet with positive
directional coefficient. Its normalized supporting functional attains
the inner minimum. The polytope is bounded and $v\ne0$, so the ray endpoint
is finite.

One explicit linear formulation of the inner dual introduces $z_{Hq}$
and $s$ with
\[
 z_{Hq}\ge\sum_{g\in q}y(g,c)\quad(H,q,c),\qquad
 s\ge\sum_qz_{Hq}\quad(H\in\mathcal H),\qquad
 \langle y,v\rangle=1,
\]
and minimizes $s-\langle y,\pi\rangle$. The primal formulation is obtained
by replacing the right side of \eqref{eq:r74-quotient-lp} by $\pi+tv$.
It is bounded and contains $t=0$. Vertex arithmetic proves the endpoint
and controller assertions, and the finite maximum over admissible families
gives the budget endpoint. These are exact linear certificates in the
specified finite problem, not a formal proof-assistant verification.
\end{proof}

\subsection{Fourier support and a sharp positivity transition}
For an abelian $G$, let $S$ be the complex characters whose vector Fourier
coefficient in $b$ is nonzero. It contains the trivial character and is
closed under conjugation.

\begin{corollary}[The abelian module cover]
\label{cor:r74-abelian-cover}
The number $\nu_G(V_b)$ is the minimum
$\sum_{H\in\mathcal H}[G:H]$ over distinct subgroup families such that
each character in $S$ is trivial on at least one $H\in\mathcal H$.
This lower bound is the exact value of $\kappa_G(b_t)$ for all sufficiently
small positive contrasts. It need not be attained at full contrast.
\end{corollary}
\begin{proof}
Character averaging projects the cyclic vector $b$ onto each of its
nonzero Fourier modes. Thus its complex cyclic module contains each
active character once; conjugate pairs give the corresponding real
two-dimensional summands. The permutation module on $H\backslash G$
contains exactly the characters trivial on $H$, each once. Containment
of $V_b$ is therefore exactly the stated cover. Repeated subgroups add no
new characters and may be deleted. Apply
Theorem~\ref{thm:r74-low-contrast}.
\end{proof}

The following example computes a complete rational threshold and an explicit
dual obstruction. Let $G=(\mathbb F_2)^4$, with the four coordinate flips as
commands, and let $\chi_1,\ldots,\chi_4$ be its coordinate sign characters.
Put $\chi_5=\chi_1\chi_2\chi_3\chi_4$ and
\begin{equation}\label{eq:r74-five-character-response}
 \beta_t(g)=\frac38+\frac t8\sum_{i=1}^5\chi_i(g),\qquad
 b_t(g)=(\beta_t(g),1-\beta_t(g)),\qquad 0\le t\le1.
\end{equation}

\begin{proposition}[A sharp transition invisible to the observable module]
\label{prop:r74-circuit-transition}
For every $t>0$, \eqref{eq:r74-five-character-response} has observable
group $(\mathbb F_2)^4$, dimension six, residual cardinality sixteen, and
$\nu_G(V_b)=10$. Its exact positive core and pure-start memory are
\[
 \kappa_G(b_t)=
 \begin{cases}10,&0<t\le3/5,\\12,&3/5<t\le1,\end{cases}
 \qquad
 m(b_t)=
 \begin{cases}11,&0<t\le3/5,\\13,&3/5<t\le1.\end{cases}
\]
At zero contrast the minimum is one. In particular a twelve-label minimum
does not occur in this family.
\end{proposition}
\begin{proof}
The five signs have product one, so their sum is $5,1$, or $-3$.
Consequently $\beta_1$ takes the values $1,1/2,0$, and every $\beta_t$
is a valid convex combination with the constant $3/8$.
Any four of the five characters are independent over $\mathbb F_2$.
For $t>0$ all five character coefficients and the constant mode are
nonzero, giving dimension six and trivial joint kernel. Thus $G$ is the
observable group and $N=16$.

For a quotient orbit write $W=H^\perp$ and $r=\dim W$; its size is $2^r$.
A proper $W$ contains at most $r$ of the five active nontrivial characters,
and $2^r\ge2r$ for $r=1,2,3$. An orbit with $r=4$ costs sixteen.
Covering the five modes therefore costs at least ten, and the five
two-point quotients attain this module bound.

Suppose that a positive core has at most eleven points. A size-eight
orbit covers at most three active characters; the remaining budget at
most three covers at most one more. Size sixteen is unavailable.
The nontrivial orbits thus have sizes two or four. Their total cost
must be ten and their active-character blocks partition the five modes
into singletons or pairs: each dimension must carry an active character
and no active character can be covered twice. At most one fixed point
can remain. Products belonging to different pairs are distinct and
inactive, because equality would give a nontrivial relation involving
at most four of the five circuit characters.

For this family of blocks let $\mathcal B_2$ denote its pairs and define
\[
 z(g)=1-\sum_{i=1}^5\chi_i(g)
       +\sum_{\{i,j\}\in\mathcal B_2}\chi_i(g)\chi_j(g).
\]
Averaging $z$ over the stabilizer of a singleton orbit gives $1-\chi_i$;
over a pair orbit it gives $(1-\chi_i)(1-\chi_j)$; over a fixed orbit
it gives $1$. The absence of any relation of length at most four ensures
that no other term survives these averages. Each resulting function is
nonnegative. Hence any nonnegative quotient contributions satisfy
\[
 \frac1{16}\sum_gz(g)\beta(g)
 =\sum_H\frac1{16}\sum_gz(g)f_H(Hg,1)\ge0.
\]
For the target \eqref{eq:r74-five-character-response}, character
orthogonality gives the value $(3-5t)/8$. Thus every core of size at
most eleven requires $t\le3/5$. This is also an explicit instance of
Theorem~\ref{thm:r74-quotient-dual}: take
$y(g,1)=-z(g)/16$, $y(g,0)=0$. Its family support value is zero and its
strict margin above the threshold is $(5t-3)/8$.

For $0<t\le3/5$, use five two-point quotient orbits, putting initial
mass $1/5$ at the plus point of each. The first-coordinate decoder on
orbit $i$ is
\[
 D_i(g,1)=\frac38+\frac{5t}{8}\chi_i(g).
\]
Its values lie in $[0,1]$, and the weighted sum is $\beta_t$. This
attains the ten-point core.

For $t=1$, use the four-point quotient on $(\chi_1,\chi_2)$ and the
eight-point quotient on $(\chi_3,\chi_4,\chi_5)$, with initial mass
$1/2$ at the all-plus point of each. Their decoders are respectively
the indicator of two plus signs and the indicator of at least two plus
signs. For signs $a,b,c$, the latter is
$1/2+(a+b+c-abc)/4$. The weighted sum of the two decoders is
\[
 \frac{1+\chi_1+\chi_2+\chi_1\chi_2}{8}
 +\frac14+
   \frac{\chi_3+\chi_4+\chi_5-\chi_3\chi_4\chi_5}{8}
 =\beta_1,
\]
because $\chi_1\chi_2=\chi_3\chi_4\chi_5$. Replace each decoder $D$ by
$(1-t)3/8+tD$ for other contrasts. This proves the twelve-point upper
everywhere, matching the strict lower above $3/5$.
The pure-start formula now gives eleven and thirteen, since $N=16$.
\end{proof}

\begin{corollary}[A robust finite-table obstruction]
\label{cor:r74-circuit-robust}
Fix $3/5<t\le1$ and a binary probability table $\widetilde\beta$ on the
same group $G=(\mathbb F_2)^4$. Supply it under the same four generator
commands, so the whole response is $w\mapsto\widetilde\beta(g(w))$.
If
\begin{equation}\label{eq:r74-circuit-robust}
 \max_{g\in G}|\widetilde\beta(g)-\beta_t(g)|
 <\frac{5t-3}{18},
\end{equation}
its exact positive core has at least twelve points and its pure-start
memory is at least thirteen. This is a sufficient lower certificate on
finite group tables. It does not assert a sharp neighbourhood radius or a
matching thirteen-label upper for every perturbed table.
\end{corollary}
\begin{proof}
Write $\eta$ for the left side of \eqref{eq:r74-circuit-robust}.
Every normalized Fourier coefficient changes by at most $\eta$.
Since $(5t-3)/18<t/8$ for $t\le1$, all five prescribed coefficients
remain nonzero. Their joint kernel is trivial, so the perturbed table
still has observable group $G$ and sixteen residuals. An action of size
at most eleven must therefore cover those same five modes, and the
preceding proof gives its singleton-and-pair partition even if additional
Fourier coefficients are present.

There are zero, one, or two pairs in such a partition. Every permutation
of the five circuit characters extends to a linear automorphism of the
dual group, because any four are a basis and the fifth is their sum.
Thus the absolute sum of its witness $z$ depends only on the number of
pairs. Direct evaluation on the sixteen sign patterns gives the following
values and multiplicities:
\[
 \begin{array}{c|l|c}
 \text{pairs}&\text{value : multiplicity of }z
       &16^{-1}\sum_g|z(g)|\\ \hline
 0&-4:1,\ 0:10,\ 4:5&3/2\\
 1&-3:1,\ -1:6,\ 1:4,\ 3:2,\ 5:3&17/8\\
 2&-2:5,\ 0:4,\ 2:2,\ 4:4,\ 6:1&9/4
 \end{array}
\]
The normalized pairing with $\widetilde\beta$ is at most
\[
 \frac{3-5t}{8}+\frac94\eta<0.
\]
It cannot have the required nonnegative quotient decomposition. Hence
$\kappa_G(\widetilde b)\ge12$, and
$m(\widetilde b)=\min\{16,1+\kappa_G(\widetilde b)\}\ge13$.
The upper thirteen on the exact contrast line was supplied by a specific
construction and is not inferred for arbitrary perturbations.
\end{proof}

For an exact physical realization one may use the sixteen-dimensional
diagonal model: prepare the identity basis state, use regular permutation
unitaries for the four generators, and take the diagonal terminal effect
with entries $\beta_t(g)$. This is already classical; no quantum advantage
is asserted for this example. The general low-contrast theorem, the normal
quotient linear programs, and this exact transition have different scopes.
